% =========================================================
\section{Referências}
% =========================================================

Embora as referências não constituam propriamente um capítulo argumentativo,
a construção da bibliografia também pode ser orientada por perguntas.

Note que, utilizando o \LaTeX, todas as referências deve estar contidas no arquivo \texttt{Referencias.bib}.Sendo, evocadas no TCC pelos comandos abaixo, dentre outros: 
\texttt{/\/cite{}, /\/citeonline{},/\/citeauthor{bibid},/\/citeauthoronline{}}.


\begin{enumerate}
	\item Todas as afirmações científicas relevantes estão devidamente referenciadas?
	\item As referências utilizadas são primárias sempre que possível?
	\item Os trabalhos mais relevantes e recentes foram considerados?
	\item Os trabalhos seminais foram identificados?
	\item As referências são provenientes de fontes academicamente confiáveis?
	\item Os artigos citados realmente sustentam as afirmações feitas?
	\item Há excesso de referências secundárias?
	\item Há referências antigas que deveriam ser complementadas por trabalhos recentes?
	\item Todas as referências citadas no texto aparecem na lista final?
	\item Todas as referências listadas são efetivamente citadas no texto?
\end{enumerate}

Veja o exemplo a seguir:

``
...
A análise da estabilidade e do comportamento transitório em malha fechada é um passo crítico no desenvolvimento de sistemas de automação industrial. Para mitigar erros em regime permanente sem comprometer a velocidade de resposta do sistema, a literatura técnica frequentemente recorre à sintonia de controladores Proporcional-Integral-Derivativo (PID). 

Conforme aponta \citeonline{ogata1998engenharia}, a ação integradora elimina o erro residual frente a pertubações em degrau, enquanto a ação derivativa antecipa a tendência do erro, adicionando amortecimento ao sistema. Desse modo, o projeto de compensadores baseados no método do lugar das raízes \cite{ogata1998engenharia} permite posicionar os polos dominantes de forma a satisfazer os requisitos de sobressinal e tempo de acomodação predefinidos.

...

''